\documentclass[12pt, a4paper]{report}

% --- UNIVERSAL PREAMBLE BLOCK ---
\usepackage[
  margin=2cm
]{geometry}
\usepackage{fontspec}

\usepackage[turkish, bidi=basic, provide=*]{babel}

\babelprovide[import, onchar=ids fonts]{turkish}
\babelprovide[import, onchar=ids fonts]{english}

\usepackage{enumitem}
\setlist[itemize]{label=-}

\usepackage{dirtree}
\usepackage{listings}
\usepackage{xcolor}
\usepackage{booktabs}
\usepackage{parskip}
\usepackage{placeins}
\usepackage{float}
\usepackage{tikz}
\usetikzlibrary{positioning, shapes.multipart, arrows.meta}

%----------BOX----------

\usepackage{tcolorbox}

\newtcolorbox{block}[1]{
  colback=blue!5!white,
  colframe=blue!75!black,
  fonttitle=\bfseries,
  title=#1,
  arc=1mm,
  sharp corners=south
}

%----------METADATA----------
\usepackage{titling}
\usepackage{titlesec}
\titleformat{\chapter}[block]{\Large\bfseries}{\thechapter.}{0.5em}{}
\titlespacing*{\chapter}{0pt}{-10pt}{8pt}

\newcommand{\projectname}{Proje 2: Veritabanı Yedekleme ve Felaketten Kurtarma Planı}
\newcommand{\studentnumber}{22290121}

\title{BLM4522 - Ağ Tabanlı Paralel Dağıtım Sistemleri Vize Projesi}
\author{Melike Vurucu}
\date{\today}

%----------TITLE LAYOUT----------
\pretitle{
\begin{center}\LARGE}
  \posttitle{\\[0.5em]\large\projectname
\end{center}}
\preauthor{
\begin{center}}
  \postauthor{\ (\studentnumber)\\[0.3em]\small Github: \url{https://github.com/melikechan/blm4522}
\end{center}}

%----------CODE STYLE----------
\definecolor{codegray}{rgb}{0.95,0.95,0.95}
\definecolor{codeblue}{rgb}{0.0,0.3,0.6}
\definecolor{codegreen}{rgb}{0.1,0.5,0.1}
\definecolor{codepurple}{rgb}{0.4,0.0,0.5}
\definecolor{codered}{rgb}{0.7,0.0,0.0}

\lstdefinestyle{sqlstyle}{
  backgroundcolor=\color{codegray},
  commentstyle=\color{codegreen}\itshape,
  keywordstyle=\color{codeblue}\bfseries,
  stringstyle=\color{codered},
  basicstyle=\ttfamily\small,
  breaklines=true,
  captionpos=b,
  keepspaces=true,
  showspaces=false,
  showstringspaces=false,
  tabsize=4,
  frame=single,
  framerule=0.5pt,
  language=SQL,
  extendedchars=true,
  morekeywords={
    ROLE, LOGIN, PASSWORD, CREATEDB, CREATEROLE,
    CONNECTION, LIMIT, SECURITY, DEFINER, ENABLE,
    FORCE, ROW, LEVEL, POLICY, USING, WITH, CHECK,
    TRIGGER, AFTER, EACH, EXECUTE, FUNCTION, RETURNS,
    LANGUAGE, plpgsql, DECLARE, BEGIN, END, RETURN,
    QUERY, RAISE, NOTICE, EXCEPTION, WHEN, THEN,
    REPLICATION, SUPERUSER, NOCREATEDB, NOCREATEROLE, NOSUPERUSER,
    TIMESTAMPTZ, SMALLINT, REFERENCES, CASCADE, UNIQUE
  }
}

\lstset{style=sqlstyle}

% hyperref must be loaded last
\usepackage{hyperref}
\hypersetup{
  colorlinks=true,
  linkcolor=black!60!black,
  urlcolor=blue!60!black,
  citecolor=blue!60!black
}

%----------DOCUMENT----------
\begin{document}

\maketitle
\tableofcontents

\chapter{Giriş}

Bu projede PostgreSQL veritabanı yönetim sistemi kullanılarak Linux ortamında bir yedekleme ve felaketten kurtarma çözümü uygulamalı olarak gösterilmiştir.

Proje kapsamında \texttt{hastane\_db} adlı kurgusal bir hastane yönetim sistemi veritabanı oluşturulup bu veritabanı üzerinde aşağıdaki bileşenler hayata geçirilmiştir:

\begin{enumerate}
  \item \textbf{Şifre tabanlı kimlik doğrulama:} \texttt{scram-sha-256} ile güvenli erişim
  \item \textbf{Üç farklı yedekleme stratejisi:} Tam (full), fark (differential) ve artık (incremental) yedekleme
  \item \textbf{Otomatik zamanlayıcı:} Cron ile yedekleme işlerinin belirli aralıklarla otomatikleştirilmesi
  \item \textbf{Felaketten kurtarma:} Kaza ile silinen, düşürülen veya yanlış güncellenen verilerin geri kazanılması
  \item \textbf{Point-in-Time Recovery (PITR):} WAL replay ile belirli bir zaman noktasına geri dönme
  \item \textbf{Yedek doğrulama:} Alınan yedeklerin gerçekten kullanılabilir olduğunun test edilmesi
\end{enumerate}

\section{Kullanılan Teknolojiler}

\begin{itemize}
  \item \textbf{PostgreSQL 18}: Açık kaynaklı, nesne-ilişkisel veritabanı yönetim sistemi.
  \item \textbf{pg\_dump / pg\_restore}: PostgreSQL'in mantıksal yedekleme araçları. Özel format (\texttt{-Fc}), maksimum sıkıştırma (\texttt{-Z9}), tek tablo geri yükleme (\texttt{--table}) desteklenmektedir.
  \item \textbf{pg\_basebackup}: Fiziksel temel yedek almak için kullanılan araç. Artık yedeklemenin başlangıç noktasını oluşturur.
  \item \textbf{WAL (Write-Ahead Logging)}: PostgreSQL'in değişiklikleri önce bir günlük dosyasına yazdığı mekanizma. WAL arşivleme etkinleştirildiğinde bu dosyalar ayrı bir dizine kopyalanır ve PITR için kullanılır.
  \item \textbf{Bash}: Kurulum ve yapılandırma adımlarını otomatikleştirmek için kullanılan kabuk (shell) betik dili.
  \item \textbf{Cron}: Yedekleme görevlerini belirli aralıklarla otomatik çalıştırmak için kullanılmıştır.
  \item \textbf{scram-sha-256}: RFC 7677 tabanlı kimlik doğrulama protokolü; PostgreSQL 14'ten itibaren varsayılan olarak kullanılmaktadır.
\end{itemize}

\newpage
\section{Veritabanı Şeması}

\textbf{hastane\_db} veritabanı bir hastane yönetim sistemini simüle eden sekiz tablodan oluşmaktadır. Şekil~\ref{fig:er}'deki entity-relationship diyagramı tablo ilişkilerini göstermektedir.

\tikzset{
  ent/.style={
    rectangle split, rectangle split parts=2,
    draw=blue!60!black, very thick,
    rectangle split part fill={blue!15!white, white},
    font=\small\ttfamily,
    align=left,
    inner sep=5pt,
    text width=3.2cm
  }
}

\begin{figure}[H]
  \centering
  \resizebox{\linewidth}{!}{%
    \begin{tikzpicture}[
        fk/.style={->, >=Stealth, dashed, thick, blue!50!black, line width=1.2pt}
      ]

      \node[ent] (bol) at (0,0) {
        \normalfont\bfseries\small\sffamily bolumler
        \nodepart{second}
        \underline{id}\ PK\\[1pt]
        ad\ UNIQUE\\[1pt]
        kat, kapasite
      };

      \node[ent] (dok) at (5.5,0) {
        \normalfont\bfseries\small\sffamily doktorlar
        \nodepart{second}
        \underline{id}\ PK\\[1pt]
        tc\_no\ UNIQUE\\[1pt]
        ad, soyad, uzmanlik\\[1pt]
        bolum\_id\ FK\\[1pt]
        email, telefon
      };

      \node[ent] (has) at (11,0) {
        \normalfont\bfseries\small\sffamily hastalar
        \nodepart{second}
        \underline{id}\ PK\\[1pt]
        tc\_no\ UNIQUE\\[1pt]
        ad, soyad\\[1pt]
        dogum\_tarihi, cinsiyet\\[1pt]
        kan\_grubu
      };

      \node[ent] (yat) at (0,-5.5) {
        \normalfont\bfseries\small\sffamily yataklar
        \nodepart{second}
        \underline{id}\ PK\\[1pt]
        bolum\_id\ FK\\[1pt]
        numara, durum
      };

      \node[ent] (ran) at (5.5,-5.5) {
        \normalfont\bfseries\small\sffamily randevular
        \nodepart{second}
        \underline{id}\ PK\\[1pt]
        hasta\_id\ FK\\[1pt]
        doktor\_id\ FK\\[1pt]
        tarih\_saat, durum
      };

      \node[ent] (yts) at (11,-5.5) {
        \normalfont\bfseries\small\sffamily yatislar
        \nodepart{second}
        \underline{id}\ PK\\[1pt]
        hasta\_id\ FK\\[1pt]
        doktor\_id\ FK\\[1pt]
        yatak\_id\ FK\\[1pt]
        giris/cikis\_tarihi
      };

      \node[ent] (ted) at (3,-10) {
        \normalfont\bfseries\small\sffamily tedaviler
        \nodepart{second}
        \underline{id}\ PK\\[1pt]
        hasta\_id\ FK\\[1pt]
        doktor\_id\ FK\\[1pt]
        tani, tarih
      };

      \node[ent] (rec) at (8.5,-10) {
        \normalfont\bfseries\small\sffamily receteler
        \nodepart{second}
        \underline{id}\ PK\\[1pt]
        tedavi\_id\ FK\\[1pt]
        ilac\_adi, doz\\[1pt]
        sure\_gun
      };

      % FK arrows
      \draw[fk] (dok.west)  -- (bol.east)   node[midway,above,font=\footnotesize\sffamily]{N:1};
      \draw[fk] (yat.north) -- (bol.south)  node[midway,right,font=\footnotesize\sffamily]{N:1};
      \draw[fk] (ran.west)  -- ++(-0.6,0) |- (dok.south) node[pos=0.2,left,font=\footnotesize\sffamily]{N:1};
      \draw[fk] (ran.east)  -- ++(0.6,0)  |- (has.south) node[pos=0.2,right,font=\footnotesize\sffamily]{N:1};
      \draw[fk] (yts.west)  -- (ran.east)  node[midway,above,font=\footnotesize\sffamily]{};
      \draw[fk] (yts.north) -- (has.south) node[midway,right,font=\footnotesize\sffamily]{N:1};
      \draw[fk] (yts.west)  -- ++(-0.4,0) |- (yat.east) node[pos=0.2,above,font=\footnotesize\sffamily]{N:1};
      \draw[fk] (ted.north) -- ++(0,0.5)  -| (has.south) node[pos=0.2,above,font=\footnotesize\sffamily]{N:1};
      \draw[fk] (ted.north) -- ++(0,0.5)  -| (dok.south) node[pos=0.2,above,font=\footnotesize\sffamily]{};
      \draw[fk] (rec.west)  -- (ted.east)  node[midway,above,font=\footnotesize\sffamily]{N:1};

    \end{tikzpicture}%
  }
  \caption{hastane\_db ER diyagramı}
  \label{fig:er}
\end{figure}

\section{Proje Senaryosu}

Projede \texttt{hastane\_db} adlı kurgusal bir hastane yönetim sistemi veritabanı kullanılmıştır. Bu veritabanı gerçek bir klinik ortamı simüle eder ve sekiz tablodan oluşur:

\begin{itemize}
  \item \textbf{bolumler:} Hastane klinik bölümleri (10 kayıt)
  \item \textbf{doktorlar:} Doktor bilgileri ve uzmanlıkları (12 kayıt)
  \item \textbf{hastalar:} Hasta kişisel ve tıbbi bilgileri (50 kayıt)
  \item \textbf{yataklar:} Bölümlere ait yatak envanteri (50 kayıt)
  \item \textbf{randevular:} Hasta-doktor randevuları (120 kayıt)
  \item \textbf{yatislar:} Yatarak tedavi kayıtları (15 kayıt)
  \item \textbf{tedaviler:} Tanı ve tedavi kayıtları (80 kayıt)
  \item \textbf{receteler:} Tedaviye bağlı ilaç reçeteleri (100+ kayıt)
\end{itemize}

\section{Proje Klasör Yapısı}

Aşağıda projenin klasör yapısı verilmiştir:

\vspace{1em}
\dirtree{%
  .1 proje2-yedekleme/.
  .2 config/.
  .3 backup.env \qquad \# Parolalar, dizin yolları (gitignore'da).
  .3 backup.env.sample \qquad \# Şablon — sürüm kontrolünde.
  .3 pg\_hba\_example.conf \qquad \# Kimlik doğrulama şablonu.
  .3 postgresql\_example.conf \qquad \# WAL yapılandırma şablonu.
  .3 pitr\_recovery\_example.conf \qquad \# PITR kurtarma örneği.
  .2 scripts/.
  .3 00\_kurulum.sh \qquad \# Kurulum (şema + veri + roller).
  .3 01\_db\_olustur.sql \qquad \# Veritabanı şeması ve view'lar.
  .3 02\_ornek\_veri\_ekle.sql \qquad \# Örnek veri (400+ kayıt).
  .3 03\_rol\_ekle.sql \qquad \# Roller ve yetkiler.
  .3 04\_tam\_yedekleme.sh \qquad \# Tam yedek (4 format).
  .3 05\_artik\_yedekleme.sh \qquad \# Artık yedek (WAL + pg\_basebackup).
  .3 06\_fark\_yedekleme.sh \qquad \# Fark yedek.
  .3 07\_zamanlanmis\_yedekleme.sh \qquad \# Cron kurulumu.
  .3 08\_felaket\_kurtarma.sh \qquad \# 4 felaket senaryosu.
  .3 09\_pitr\_kurtarma.sh \qquad \# Point-in-Time Recovery.
  .3 10\_yedek\_test.sh \qquad \# Yedek doğrulama ve testler.
}

\chapter{Yedekleme Stratejileri}

\texttt{00\_kurulum.sh} dört adımda çalışır:

\begin{enumerate}
  \item yedekleme dizinlerini oluşturur,
  \item \texttt{01\_db\_olustur.sql} ile şemayı oluşturur,
  \item \texttt{02\_ornek\_veri\_ekle.sql} ile örnek verileri yükler,
  \item \texttt{03\_rol\_ekle.sql} ile rolleri ve yetkileri yapılandırır.
\end{enumerate}

\texttt{config/postgresql\_example.conf} dosyasındaki WAL yapılandırması \texttt{postgresql.conf}'a eklenir:

\begin{lstlisting}[caption={WAL arşivleme yapılandırması (postgresql\_example.conf)}]
-- WAL seviyesi: PITR ve replikasyon icin gerekli
wal_level = replica

-- WAL dosyalarini arsiv dizinine kopyala
archive_mode = on
archive_command = 'cp %p /var/lib/postgres/backups/wal_archive/%f'

-- pg_basebackup icin gonderici sayisi
max_wal_senders = 3
\end{lstlisting}

\section{Kimlik Doğrulama Ayarları}

Varsayılan Arch Linux kurulumunda yerel bağlantılar için \texttt{peer} kullanılır. Bu projede \texttt{scram-sha-256} etkinleştirilmiştir:

\begin{lstlisting}[caption={pg\_hba\_example.conf — şifre doğrulama}]
-- postgresql.conf
password_encryption = scram-sha-256

-- pg_hba.conf
local   hastane_db   backup_user                    scram-sha-256
local   hastane_db   hastane_app                    scram-sha-256
host    hastane_db   backup_user   127.0.0.1/32     scram-sha-256
host    hastane_db   hastane_app   127.0.0.1/32     scram-sha-256
-- WAL replikasyonu
local   replication  backup_user                    scram-sha-256
host    replication  backup_user   127.0.0.1/32     scram-sha-256
\end{lstlisting}

\newpage
İki ayrı rol oluşturulmuştur; yedekleme araçlarının yetkisi uygulama kullanıcısından farklıdır:

\begin{lstlisting}[caption={03\_rol\_ekle.sql — roller}]
-- Rol varsa guncelle, yoksa olustur (idempotent)
DO $$
BEGIN
    IF NOT EXISTS (SELECT FROM pg_roles WHERE rolname = 'backup_user') THEN
        CREATE ROLE backup_user WITH LOGIN PASSWORD 'Backup@Hastane2025!'
            SUPERUSER REPLICATION;
    ELSE
        ALTER ROLE backup_user WITH LOGIN PASSWORD 'Backup@Hastane2025!'
            SUPERUSER REPLICATION;
    END IF;
END $$;

DO $$
BEGIN
    IF NOT EXISTS (SELECT FROM pg_roles WHERE rolname = 'hastane_app') THEN
        CREATE ROLE hastane_app WITH LOGIN PASSWORD 'HastaneApp@2025!'
            NOCREATEDB NOCREATEROLE NOSUPERUSER;
    ELSE
        ALTER ROLE hastane_app WITH LOGIN PASSWORD 'HastaneApp@2025!'
            NOCREATEDB NOCREATEROLE NOSUPERUSER;
    END IF;
END $$;

ALTER DATABASE hastane_db OWNER TO backup_user;
\end{lstlisting}

\newpage
\section{Tam Yedekleme (Full Backup)}

Tam yedekleme, veritabanının tamamını tek bir dosyaya aktarır. Diğer yedekleme türlerinin \textbf{referans noktası}dır. \texttt{04\_tam\_yedekleme.sh} dört farklı formatta yedek alır:

\begin{lstlisting}[caption={04\_tam\_yedekleme.sh — dört format}]
PGDUMP="pg_dump -h $DB_HOST -p $DB_PORT -U $DB_USER $DB_NAME"

-- Custom format: pg_restore uyumlu, secici geri yukleme destekler
$PGDUMP -Fc -Z9 -f "${DB_NAME}_tam_${TS}.dump"

-- SQL format: her PostgreSQL surumuyle uyumlu
$PGDUMP -Fp -f "${DB_NAME}_tam_${TS}.sql"

-- Gzip SQL: maksimum disk tasarrufu
$PGDUMP -Fp | gzip -9 > "${DB_NAME}_tam_${TS}.sql.gz"

-- Tar arsivi: dizin tabanli geri yukleme
$PGDUMP -Ft -f "${DB_NAME}_tam_${TS}.tar"

-- Saklama politikasi: eski yedekleri sil
find "$FULL_DIR" -type f -mtime "+${RETENTION_DAYS}" -delete
\end{lstlisting}

\begin{table}[H]
  \centering
  \caption{Yedekleme formatları ve özellikleri}
  \begin{tabular}{llll}
    \toprule
    \textbf{Format} & \textbf{Uzantı} & \textbf{Geri Yükleme} & \textbf{Özellik} \\
    \midrule
    Custom   & \texttt{.dump}   & \texttt{pg\_restore}    & Seçici geri yükleme \\
    SQL      & \texttt{.sql}    & \texttt{psql}           & İnsan okunabilir \\
    Gzip SQL & \texttt{.sql.gz} & \texttt{zcat | psql}    & En küçük boyut \\
    Tar      & \texttt{.tar}    & \texttt{pg\_restore}    & Dizin bazlı \\
    \bottomrule
  \end{tabular}
  \label{tab:backup-formats}
\end{table}

\newpage
\section{Artık Yedekleme (Incremental Backup)}

Artık yedekleme, \textbf{son herhangi bir yedekten} bu yana üretilen WAL dosyalarını arşivler. \texttt{05\_artik\_yedekleme.sh} temel yedek alır ve WAL dosyalarını kopyalar:

\begin{lstlisting}[caption={05\_artik\_yedekleme.sh — WAL arşivleme}]
-- Temel yedek: pg_basebackup ile fiziksel kopya
pg_basebackup \
    -h "$DB_HOST" -p "$DB_PORT" -U "$DB_USER" \
    -D "$BASE_DIR" -Ft -z --checkpoint=fast -P

-- Aktif WAL segmentini tamamla
psql ... -c "SELECT pg_switch_wal();"

-- WAL dosyalarini postgres yetkisiyle oku, hedef dizine yaz
while IFS= read -r wal_file; do
    dest="$WAL_DIR/$(basename "$wal_file")"
    if [[ ! -f "$dest" ]]; then
        sudo -u postgres cat "$wal_file" > "$dest"
    fi
done < <(sudo -u postgres find "$PG_DATA_DIR/pg_wal" \
    -maxdepth 1 -type f -name '[0-9A-F]*')
\end{lstlisting}

\newpage
\section{Fark Yedekleme (Differential Backup)}

Fark yedekleme, \textbf{her zaman son TAM yedekten} bu yana değişen verileri arşivler. \texttt{06\_fark\_yedekleme.sh} dört farklı çıktı üretir:

\begin{lstlisting}[caption={06\_fark\_yedekleme.sh — fark yedek türleri}]
-- 1. Sema-only: tablolar, indeksler, view'lar
$PGDUMP --schema-only -f "sema_only_${TS}.sql"

-- 2. Veri-only: yalnizca INSERT ifadeleri
$PGDUMP --data-only --column-inserts -f "veri_only_${TS}.sql"

-- 3. Secici tablo yedekleri (sik degisen tablolar)
for tablo in randevular yatislar tedaviler receteler; do
    $PGDUMP -Fc -t "$tablo" -f "${tablo}_${TS}.dump"
done

-- 4. Son 7 gun degisen kayitlar (zaman damgasi filtreli)
SELECT ... FROM randevular WHERE created_at >= NOW() - INTERVAL '7 days';
SELECT ... FROM tedaviler  WHERE tarih      >= NOW() - INTERVAL '7 days';
SELECT ... FROM yatislar   WHERE giris_tarihi >= NOW() - INTERVAL '7 days';
\end{lstlisting}

\begin{table}[H]
  \centering
  \caption{Yedekleme türleri karşılaştırması}
  \begin{tabular}{llll}
    \toprule
    \textbf{Tür} & \textbf{Referans} & \textbf{Boyut} & \textbf{Geri Yükleme} \\
    \midrule
    Tam (Full)  & —                 & Büyük  & Doğrudan \\
    Fark (Diff) & Son TAM yedek     & Orta   & Full + Diff \\
    Artık (Incr)& Son herhangi yedek& Küçük  & Full + Incr\textsubscript{1..n} \\
    \bottomrule
  \end{tabular}
  \label{tab:backup-types}
\end{table}

\chapter{Kurtarma ve Otomasyon}

\section{Felaketten Kurtarma Senaryoları}

\texttt{08\_felaket\_kurtarma.sh}, dört farklı felaket senaryosunu sırasıyla simüle eder ve her birini otomatik olarak kurtarır.

\subsection{Senaryo 1: Kritik Tablo Silindi (DROP TABLE)}

\textbf{Felaket:} \texttt{tedaviler} ve \texttt{receteler} tabloları yanlışlıkla düşürülür.

\begin{lstlisting}[caption={Senaryo 1 — tablo silme simülasyonu}]
DROP TABLE receteler CASCADE;
DROP TABLE tedaviler CASCADE;
\end{lstlisting}

\textbf{Kurtarma:} \texttt{pg\_restore --table} ile yalnızca silinen tablolar seçici olarak geri yüklenir:

\begin{lstlisting}[caption={Senaryo 1 — tablo geri yükleme}]
-- Sema geri yukle
pg_restore -d "$DB_NAME" --table=tedaviler --schema-only "$BACKUP"
pg_restore -d "$DB_NAME" --table=receteler --schema-only "$BACKUP"

-- Veriyi geri yukle
pg_restore -d "$DB_NAME" \
    --table=tedaviler --table=receteler \
    --data-only "$BACKUP"
\end{lstlisting}

\subsection{Senaryo 2: Toplu Veri Silindi (DELETE)}

\textbf{Felaket:} Tüm hasta kayıtları ve bağlı tablolar yanlışlıkla silinir.

\begin{lstlisting}[caption={Senaryo 2 — toplu silme}]
DELETE FROM randevular;
DELETE FROM yatislar;
DELETE FROM tedaviler;
DELETE FROM hastalar;
\end{lstlisting}

\textbf{Kurtarma:} Yedek geçici bir veritabanına geri yüklenir; \texttt{COPY} ile üretim veritabanına aktarılır:

\begin{lstlisting}[caption={Senaryo 2 — COPY ile veri kurtarma}]
-- Gecici DB olustur ve yedegi yukle
psql ... -c "CREATE DATABASE ${DB_NAME}_tmp OWNER $DB_USER;"
pg_restore ... -d "${DB_NAME}_tmp" "$BACKUP"

-- Her tabloyu uretim DB'sine kopyala
for tablo in hastalar yatislar randevular tedaviler; do
    psql -d "${DB_NAME}_tmp" -c "\COPY $tablo TO STDOUT" | \
    psql -d "$DB_NAME"       -c "\COPY $tablo FROM STDIN"
done
\end{lstlisting}

\subsection{Senaryo 3: Yanlış Toplu Güncelleme (UPDATE)}

\textbf{Felaket:} Tüm randevular yanlışlıkla \texttt{iptal} durumuna getirilir.

\begin{lstlisting}[caption={Senaryo 3 — yanlış güncelleme ve kurtarma}]
-- Felaket
UPDATE randevular SET durum = 'iptal';

-- Kurtarma: yedekten dogru durumlari al
CREATE TEMP TABLE randevu_kurtarma (id int, durum text);
\COPY randevu_kurtarma FROM '/tmp/randevu_kurtarma.csv' CSV;

UPDATE randevular r
SET durum = k.durum
FROM randevu_kurtarma k
WHERE r.id = k.id;
\end{lstlisting}

\subsection{Senaryo 4: Veritabanı Tamamen Kayboldu}

\textbf{Felaket:} Veritabanının tamamı kaybolur.

\textbf{Kurtarma:} Tam yedek yeni bir veritabanına geri yüklenir:

\begin{lstlisting}[caption={Senaryo 4 — tam kurtarma}]
-- Kurtarma veritabani olustur
psql ... -c "CREATE DATABASE ${DB_NAME}_kurtarildi OWNER $DB_USER;"

-- Tam yedegi geri yukle
pg_restore ... -d "${DB_NAME}_kurtarildi" "$BACKUP_S4"

-- Uretimle degistirmek icin:
-- ALTER DATABASE hastane_db RENAME TO hastane_db_eski;
-- ALTER DATABASE hastane_db_kurtarildi RENAME TO hastane_db;
\end{lstlisting}

\section{Point-in-Time Recovery (PITR)}

\texttt{09\_pitr\_kurtarma.sh}, WAL replay mekanizmasını kullanarak veritabanını belirli bir zaman noktasına geri döndürür.

\begin{lstlisting}[caption={PITR — geri yükleme noktası oluşturma}]
-- 1. Geri yukleme noktasi olustur
SELECT pg_create_restore_point('pitr_test_20260421');

-- 2. Temel yedek al
pg_dump -Fc -Z9 "$DB_NAME" -f "pitr_temel.dump"

-- 3. Felaket olustur — FK sirasina gore bagimli tablolar once silinir
WITH hedef AS (SELECT id FROM hastalar WHERE id % 5 = 0)
DELETE FROM receteler
    WHERE tedavi_id IN (SELECT id FROM tedaviler
                        WHERE hasta_id IN (SELECT id FROM hedef));
WITH hedef AS (SELECT id FROM hastalar WHERE id % 5 = 0)
DELETE FROM tedaviler WHERE hasta_id IN (SELECT id FROM hedef);
WITH hedef AS (SELECT id FROM hastalar WHERE id % 5 = 0)
DELETE FROM yatislar  WHERE hasta_id IN (SELECT id FROM hedef);
WITH hedef AS (SELECT id FROM hastalar WHERE id % 5 = 0)
DELETE FROM randevular WHERE hasta_id IN (SELECT id FROM hedef);
DELETE FROM hastalar WHERE id % 5 = 0;
UPDATE doktorlar SET email = 'BOZUK_' || email WHERE id % 2 = 0;
\end{lstlisting}

\begin{lstlisting}[caption={pitr\_recovery\_example.conf — kurtarma yapılandırması}]
-- WAL arsivinden geri yukleme komutu
restore_command = 'cp /var/lib/postgres/backups/wal_archive/%f %p'

-- Geri yukleme noktasina don (SECENEK A)
recovery_target_name = 'pitr_test_20260421'
recovery_target_action = 'promote'

-- Zaman bazli alternatif (SECENEK B)
-- recovery_target_time = '2026-04-21 12:00:00 +03'
\end{lstlisting}

Gerçek PITR adımları:

\begin{enumerate}
  \item PostgreSQL'i durdur: \texttt{sudo systemctl stop postgresql}
  \item \texttt{pg\_basebackup} çıktısını veri dizinine çıkart
  \item \texttt{pitr\_recovery\_example.conf} ayarlarını \texttt{postgresql.conf}'a ekle
  \item Recovery sinyali oluştur: \texttt{sudo touch /var/lib/postgres/data/recovery.signal}
  \item PostgreSQL'i başlat; kurtarma tamamlandığında \texttt{recovery.signal} otomatik silinir
\end{enumerate}

\section{Otomatik Zamanlayıcı (Cron)}

\texttt{07\_zamanlanmis\_yedekleme.sh}, yedekleme görevlerini kullanıcının crontab'ına ekler ve otomatik çalışacak bir betik (\texttt{auto\_yedekleme.sh}) oluşturur.

\begin{lstlisting}[caption={config/crontab\_yedekleme.txt — yedekleme takvimi}]
-- Her gece 02:00 -- Gunluk tam yedek
0 2 * * *   auto_yedekleme.sh

-- Her Pazar 03:00 -- Haftalik tam yedek (tum formatlar)
0 3 * * 0   04_tam_yedekleme.sh

-- Is gunleri 09:00-18:00 arasi her saat -- Fark yedek
0 9-18 * * 1-5  06_fark_yedekleme.sh

-- Her ayin 1'i 04:00 -- Aylik artik yedek
0 4 1 * *   05_artik_yedekleme.sh
\end{lstlisting}

\begin{table}[ht]
  \centering
  \begin{tabular}{lcc}
    \toprule
    \textbf{Periyot} & \textbf{Saat} & \textbf{Tür} \\
    \midrule
    Her gece         & 02:00         & Otomatik tam yedek \\
    Her Pazar        & 03:00         & Tam yedek (4 format) \\
    İş günleri       & 09:00–18:00   & Fark yedek (her saat) \\
    Her ayın 1'i     & 04:00         & Artık yedek \\
    \bottomrule
  \end{tabular}
  \caption{Haftalık yedekleme takvimi}
\end{table}

\section{Yedek Doğrulama}

Bir yedeğin var olması, geri yüklenebilir olduğu anlamına gelmez. \texttt{10\_yedek\_test.sh} beş aşamalı doğrulama yapar:

\begin{enumerate}
  \item \textbf{Yedek oluşturma:} Custom ve SQL formatlarında yedek alınır.
  \item \textbf{Bütünlük kontrolü:} \texttt{pg\_restore --list} ile format geçerli mi, bozulma var mı?
  \item \textbf{Geri yükleme testi:} Geçici veritabanına geri yüklenir; tüm tabloların kayıt sayısı üretimle karşılaştırılır.
  \item \textbf{Veri bütünlüğü:} Yabancı anahtar kısıtlamaları, view'lar ve indeksler doğrulanır.
  \item \textbf{Performans:} Yedek süresi ve boyutu ölçülür.
\end{enumerate}

\begin{lstlisting}[caption={10\_yedek\_test.sh — bütünlük ve geri yükleme testi}]
-- Butunluk kontrolu
pg_restore --list "$F_CUSTOM" > restore_list.txt
TABLO_SAYISI=$(grep -c "TABLE DATA" restore_list.txt)

-- Gecici DB'ye geri yukle
psql ... -c "CREATE DATABASE $TEST_DB OWNER $DB_USER;"
pg_restore ... -d "$TEST_DB" "$F_CUSTOM"

-- Kayit sayilarini karsilastir
for tablo in bolumler doktorlar hastalar yataklar
             randevular yatislar tedaviler receteler; do
    PROD=$(psql -d "$DB_NAME" -tAc "SELECT COUNT(*) FROM $tablo;")
    TEST=$(psql -d "$TEST_DB" -tAc "SELECT COUNT(*) FROM $tablo;")
    -- Esit degilse FAIL raporla
done

-- FK kisitlama dogrulama
SELECT COUNT(*) FROM randevular r
WHERE NOT EXISTS (SELECT 1 FROM hastalar h WHERE h.id = r.hasta_id);
-- 0 olmasi gerekir
\end{lstlisting}

Script \texttt{exit 0} (başarı) veya \texttt{exit 1} (hata) döner; cron ile entegre edildiğinde başarısız doğrulamalar \texttt{MAILTO} ile bildirilebilir.

\begin{block}{Not}
  Yedek doğrulama adımı atlanmamalıdır. Disk bozulması veya eksik WAL dosyası yedeği sessizce kullanılamaz hale getirebilir; bu durum ancak gerçek bir felaket anında fark edilir.
\end{block}

\chapter{Sonuç}

Bu projede PostgreSQL için eksiksiz bir yedekleme ve felaketten kurtarma çözümü geliştirilmiştir. Toplam \textbf{11 script} (3 SQL, 8 bash) numaralı, düz bir dizin yapısında organize edilmiştir:

\begin{itemize}
  \item \textbf{Güvenlik}: \texttt{scram-sha-256} kimlik doğrulama, iki ayrı rol (yedekleme ve uygulama), parolaların \texttt{.gitignore} ile sürüm kontrolü dışında tutulması.
  \item \textbf{Yedekleme}: Tam (4 format), fark ve artık yedekleme; WAL arşivleme ve \texttt{pg\_basebackup} ile fiziksel yedek; saklama politikası (30 gün). Artık ve fark yedekleme WAL arşivleme olmaksızın mümkün değildir; \texttt{wal\_level = replica} ve \texttt{archive\_mode = on} zorunludur.
  \item \textbf{Kurtarma}: Dört farklı felaket senaryosu (DROP TABLE, toplu DELETE, yanlış UPDATE, tam kayıp). Seçici \texttt{pg\_restore --table} ile tablo düzeyinde geri yükleme; \texttt{COPY} tabanlı hızlı veri aktarımı.
  \item \textbf{PITR}: WAL replay ile belirli bir zaman noktasına veya isimlendirilmiş geri yükleme noktasına kurtarma.
  \item \textbf{Doğrulama}: Beş aşamalı yedek doğrulama (oluşturma, bütünlük, geri yükleme, FK kontrolü, performans).
  \item \textbf{Otomasyon}: Cron zamanlayıcısı ile günlük tam, saatlik fark ve aylık artık yedekleme.
\end{itemize}

Bu bileşenlerin birlikte çalışması, tam teşekküllü bir felaket kurtarma modeli ortaya koymaktadır: tek bir yedekleme türüne bağımlılık ortadan kalkar ve veri kaybı penceresi en aza indirilir.

\end{document}
